\documentclass[11pt]{article}

\usepackage[utf8]{inputenc}
\usepackage[T1]{fontenc}
\usepackage{geometry}
\usepackage{amsmath}
\usepackage{amssymb}
\usepackage{booktabs}
\usepackage{hyperref}
\usepackage{xurl}
\usepackage{xcolor}
\usepackage{listings}
\usepackage{enumitem}
\usepackage{parskip}
\usepackage{graphicx}
\graphicspath{{../images/}}

\geometry{margin=1in}

\lstset{
  language=Java,
  basicstyle=\ttfamily\small,
  keywordstyle=\color{blue},
  commentstyle=\color{gray},
  stringstyle=\color{teal},
  showstringspaces=false,
  columns=fullflexible,
  frame=single,
  framerule=0.2pt,
  breaklines=true
}

\setlist[itemize]{topsep=0.3em,itemsep=0.2em}
\setlist[enumerate]{topsep=0.3em,itemsep=0.2em}

% Simple per-section source line.
\newcommand{\notesources}[1]{\par\noindent{\small\color{gray}\textit{Sources:} \sloppy #1\par}}

% Unit-wise source strings for this subject (used by slides/notes).
% Path is relative to the lecture `latex/` folder (the compilation CWD).
% OOPS with Java (CSEG1044) — unit-wise sources
%
% These are short, student-facing reference strings intended to be used with:
%   - \slidesources{...}  (in beamer slides)
%   - \notesources{...}   (in lecture notes)
%
% Style inspiration: other subjects in My-Subjects/ use semicolon-separated sources
% and include an "accessed" date for web documentation.

\newcommand{\OOPSUnitISources}{Schildt, \textit{Java: A Beginner's Guide} (10e), McGraw-Hill (Java fundamentals + OOP basics).; Oracle Java Tutorials: Getting Started + Language Basics + Classes and Objects (accessed \today).; Java SE API docs (e.g., \texttt{String}, \texttt{Scanner}) (accessed \today).}

\newcommand{\OOPSUnitIISources}{Schildt, \textit{Java: The Complete Reference} (12e), McGraw-Hill (classes, inheritance, packages, interfaces).; Bloch, \textit{Effective Java} (3e), Addison-Wesley, 2018 (design choices; interfaces vs abstract classes).; Oracle Java Tutorials: Classes and Objects + Interfaces + Packages (accessed \today).; Java SE API docs (e.g., \texttt{Comparable}, \texttt{Runnable}) (accessed \today).}

\newcommand{\OOPSUnitIIISources}{Schildt, \textit{Java: The Complete Reference} (12e) (nested/inner classes, exceptions, threads, I/O).; Oracle Java Tutorials: Nested Classes + Exceptions + Concurrency + Basic I/O (accessed \today).; Java SE API docs (e.g., \texttt{Thread}, \texttt{AutoCloseable}) (accessed \today).}

\newcommand{\OOPSUnitIVSources}{Schildt, \textit{Java: The Complete Reference} (12e) (generics, lambdas, Swing, JDBC).; Bloch, \textit{Effective Java} (3e) (generics + lambdas).; Oracle Java Tutorials: Generics + Lambda Expressions + Swing + JDBC Basics (accessed \today).; Java SE API docs (e.g., \texttt{java.util.function}, \texttt{javax.swing}, \texttt{java.sql}) (accessed \today).}

\newcommand{\OOPSUnitVSources}{Schildt, \textit{Java: The Complete Reference} (12e) (Collections Framework + wrapper classes).; Oracle Java docs: Collections Framework overview (accessed \today).; Java SE API docs (\texttt{java.util} collections; wrapper classes in \texttt{java.lang}) (accessed \today).}

\newcommand{\OOPSUnitVISources}{Git SCM: \textit{Pro Git} (2e) (version control workflow), accessed \today.; GitHub Docs (repositories, issues, pull requests), accessed \today.; Oracle Java Tutorials: Swing + JDBC (accessed \today).; JUnit 5 User Guide (accessed \today).; Apache Maven Project: Getting Started (accessed \today).; Gradle User Manual: Getting Started (accessed \today).; UML class diagrams: UML 2.x overview/spec (accessed \today).}

\newcommand{\OOPSMiscWhyInterfacesSources}{Schildt, \textit{Java: The Complete Reference} (12e) (interfaces).; Bloch, \textit{Effective Java} (3e) (prefer interfaces where appropriate).; Oracle Java Tutorials: Interfaces (accessed \today).; Java SE API docs (e.g., \texttt{Comparable}, \texttt{Runnable}) (accessed \today).}



\title{Object Oriented Programming Lab (CSEG1044)\\Experiment-07\&08: Classes/Objects + Overloading (Student Handout --- Before Submission)}
\author{Tofik Ali}
\date{\today}

\begin{document}
\maketitle
\tableofcontents

\section{About this handout}
This handout gives the implementation requirements and submission checklist.
Do not copy solutions from outside sources; write and test your own code.

\section{Experiment 07 (Syllabus)}
\begin{itemize}[leftmargin=*]
  \item \textbf{Objective:} To understand classes, objects, and constructors.
  \item \textbf{Problem Statement:} Design classes with data members and member functions.
  \item \textbf{Expected Outcome:} Students will understand object creation and encapsulation.
\end{itemize}

\section{Experiment 08 (Syllabus)}
\begin{itemize}[leftmargin=*]
  \item \textbf{Objective:} To demonstrate compile-time polymorphism.
  \item \textbf{Problem Statement:} Implement method and constructor overloading in Java.
  \item \textbf{Expected Outcome:} Students will understand method resolution at compile time.
\end{itemize}

\section{What you must do (minimum requirements)}

\subsection{Task A --- Student class with validation (Exp-07)}
Create \texttt{StudentDemo.java} with a \texttt{Student} class.

\noindent Fields (private):
\begin{itemize}[leftmargin=*]
  \item \texttt{rollNo} (must be positive)
  \item \texttt{name} (must be non-empty)
  \item \texttt{cgpa} (0 to 10)
\end{itemize}

\noindent Requirements:
\begin{itemize}[leftmargin=*]
  \item Validate in constructor.
  \item Provide a safe update method \texttt{updateCgpa(double)} with validation.
  \item Override \texttt{toString()} to print object details.
\end{itemize}

\subsection{Task B --- Overloading demo (Exp-08)}
Create \texttt{OverloadingDemo.java}:
\begin{enumerate}[leftmargin=*]
  \item Create class \texttt{Box} with overloaded constructors:
  \begin{itemize}[leftmargin=*]
    \item default: unit box
    \item 1 parameter: cube
    \item 3 parameters: \texttt{w, h, d}
  \end{itemize}
  \item Add method \texttt{volume()}.
  \item Write overloaded methods:
  \begin{itemize}[leftmargin=*]
    \item \texttt{area(double r)} $\rightarrow$ circle area
    \item \texttt{area(double w, double h)} $\rightarrow$ rectangle area
  \end{itemize}
  \item In \texttt{main}, create objects using different constructors and print volumes and areas.
\end{enumerate}

\section{Hints (read only if stuck)}
\begin{itemize}[leftmargin=*]
  \item Keep fields private and expose behavior through methods.
  \item Constructor overloading can reduce duplicate initialization code using \texttt{this(...)} chaining.
  \item Overloading cannot be done using only return type; parameter list must change.
\end{itemize}

\section{What to submit (MANDATORY)}
Submit a single zipped folder containing:
\begin{itemize}[leftmargin=*]
  \item \texttt{StudentDemo.java}
  \item \texttt{OverloadingDemo.java}
  \item sample output for both programs
  \item observations (5--10 lines): why encapsulation helps; how Java chooses overloaded methods
\end{itemize}

\notesources{\OOPSUnitIISources}
\end{document}
